\section{模型假设}
\label{sec:assumptions}

本节逐条列出假设及其依据与代价，并汇总说明它们对结论的影响；每条假设在后文都有对应的检验或敏感性分析。

\subsection{几何与守恒层面的假设}

\begin{assumptionlist}
\item[\textbf{A1}] \textbf{一维轴对称：温度与含水率只沿半径方向变化，轴向均匀、端面效应忽略}。\emph{依据}：药材为圆柱、热风沿侧面均匀送风；题目只要求``到药材中心的距离''的分布。长度与直径之比 $L/(2\Rzero)=25/4=6.25$，两端面积与侧面积之比为 $R_0/L=0.08$。\emph{代价}：端面的散热与失水被平摊进径向模型，表面通量可能被高估；题目未给端面边界条件，该误差无法用题目数据约束
\item[\textbf{A2}] \textbf{周向均匀：无 $\theta$ 依赖}。\emph{依据}：圆柱对称、无周向温度梯度。\emph{代价}：若实际烘房存在风向偏置，会低估局部干燥速率
\item[\textbf{A3}] \textbf{局部热平衡：同一位置处固体骨架与水分取同一温度，$\rho\cp$ 为等效值}。\emph{依据}：采用单温度有效模型，避免引入骨架与水分之间的换热系数（题目未给出）。\emph{代价}：升温初期骨架与水分温差最大，单温度模型在该时段误差最大；本问升温速率不大，误差有限
\item[\textbf{A4}] \textbf{热方程采用\textbf{有效显热容量}闭合 $\rho\cp{\T}_\t=(\rpos\kcond{\T}_\rpos)_\rpos/\rpos$}。\emph{依据}：题目提供的物性只有 $\rho$、$\cp$、$\kcond$，没有焓输运与相变项的参数。\emph{代价}：本文\textbf{不声称}该方程是包含蒸发、骨架与水焓输运的完整热力学守恒式；能量审计只验证离散格式的瞬时收支闭合
\item[\textbf{A5}] \textbf{表面热边界为\textbf{纯对流}，不另加蒸发潜热项}。\emph{依据}：题目附录只给出对流换热系数 $\hconv$，未给出水的汽化潜热与蒸发速率的关系式。\emph{代价}：蒸发冷却缺项会使表面温度偏高、干燥偏快，在 $2\sim3$ 天的长时间干燥中影响明显大于 $30$~min 的预热阶段；本文未引入该参数（不凭空补一个题目未给的量），并在第~\ref{sec:limitations} 节列为最大的物理闭合风险
\item[\textbf{A6}] \textbf{不发生化学反应与体积热源，忽略辐射换热}。\emph{依据}：题目未给出辐射与反应参数；烘房温度不超过 $50.2~^\circ$C，辐射量级小。\emph{代价}：高温工况（$>100~^\circ$C）下该假设不再适用
\end{assumptionlist}

\subsection{边界条件与物性闭合层面的假设}

\begin{assumptionlist}
\item[\textbf{A7}] \textbf{表面传质采用题给的浓度差驱动形式 $\Ddiff{\C}_\rpos|_{\Rzero}=-\hmass({\C}_s-\Cinf)$}。\emph{依据}：题目直接给出对流传质系数 $\hmass=8\times10^{-7}$~m/s 与烘房水分浓度 $\Cinf$，形式由题面确定。\emph{代价}：\textbf{属有效闭合假设而非已证明的物理关系}：题目未说明气相水分的 kg/kg 与药材干基含水率的 kg/kg 是否同一基准（二者相差约 $51$ 倍）。若基准不一致，边界驱动势会整数量级地偏，进而使烘干时长偏。本文按题给形式使用，并在第~\ref{sec:limitations} 节明确披露
\item[\textbf{A8}] \textbf{水分控制方程按\textbf{单位干物质质量}书写，干基表观密度 $\rho_d$ 从方程两端消去；经验公式给出的含湿表观密度 $\rho$ 只进入热容量 $\rho\cp$}。\emph{依据}：干基含水率 $\C$ 的定义就是水质量与干物质质量之比；$\rho_d$ 未经题目给出，若强行将其代入会凭空引入一个参数。\emph{代价}：质量审计是在\textbf{材料测度}（$\sum_i V_i{\C}_i$）下做的，不是物理体积口径；$\rho$ 随 $\C$ 变化对水分输运的间接影响被忽略
\item[\textbf{A9}] \textbf{附录 2 的参数集下温度方程与水分方程解耦（$\Ddiff$ 不含 $\T$）；附录 3/4 下二者双向弱耦合}。\emph{依据}：由题目给出的公式形式直接判断。\emph{代价}：无额外代价；但耦合使数值求解的非线性增强，需要隐式方法处理
\item[\textbf{A10}] \textbf{问题 4 中\textbf{长度不变、半径按附件 2 收缩}，且收缩在径向是\textbf{均匀}的（同一物质壳层按同一比例缩放）}。\emph{依据}：方程按单位长度写出，长度不进入任何一项；附件 2 只给出半径，未给长度与体积。\emph{代价}：这是问题 4 最大的几何不确定性：若长度按同一比例缩短，药材体积将只剩 $21.5\%$（而非固定长度假设下的 $35.9\%$）；附件 2 的数据无法区分这两种情形。此外不允许出现``外干内湿''的非均匀收缩
\item[\textbf{A11}] \textbf{问题 4 中\textbf{不引入吸附平衡含水率 ${\C}_e$}，$\Cinf$ 只作为边界驱动势}。\emph{依据}：题目只给出终点判据 $\C<0.15$~kg/kg，未给吸附等温线；引入 ${\C}_e$ 等于凭空增加一个题目未给出的参数。\emph{代价}：本文不声称存在使干燥停滞的正平衡含水率；末期表面含水率渐近趋近但不等于 $\Cinf$，这是``$\Cinf$ 仅作驱动势''的直接后果
\end{assumptionlist}

\subsection{数据与判据层面的假设}

\begin{assumptionlist}
\item[\textbf{B1}] \textbf{附件 1 的测量时刻之间用\textbf{分段线性插值}}。\emph{依据}：题目只给离散测量值；线性插值不会在相邻数据之间制造额外极值，且不引入求解器难以处理的振荡。\emph{代价}：插值函数的一阶导数在数据点处不连续；本文在每个附件测量节点处重启积分器，避免自适应步长跨越折点
\item[\textbf{B2}] \textbf{附件 1 只覆盖 $0\sim4$~h，\textbf{$t>4$~h 后环境条件按末值保持外推}（$\Tinf=50.165~^\circ$C，$\Cinf=0.04986$~kg/kg）}。\emph{依据}：题目说``预热平衡与恒温干燥阶段的参数有所不同''，暗示恒温段环境稳定；题目未给出 $4$~h 之后的环境数据。\emph{代价}：\textbf{这是凭空延伸数据，是本文最易被质疑之处}：该外推覆盖问题 2/3/4 全程约 $93\%$ 的时长。本文在第~\ref{sec:sensitivity} 节按附件 1 末小时的实测上下限做灵敏度分析，量化其对烘干时长的影响（小时量级）
\item[\textbf{B3}] \textbf{问题 3、4 的终点判据取\textbf{全场最大值}低于 $0.15$~kg/kg（见式~\eqref{eq:maxcrit}）}。\emph{依据}：水分由内向外扩散，轴心最慢达标，故最大值判据是最保守的解释。\emph{代价}：若按断面平均值理解，烘干时间会更短；判据取值依赖内部网格（网格越细最大值越大），本文用网格加密给出其影响量级
\item[\textbf{B4}] \textbf{判据在\textbf{未舍入}的内部解上判断；交付文件的四位小数只是输出格式}。\emph{依据}：四位小数是题目的展示要求，不是物理精度；若用舍入后的交付表判断达标，会出现相邻两行都显示 $0.1500$ 而无法区分的情形。\emph{代价}：报出的烘干时长精度受积分器容差与网格限制，本文用时间方法交叉验证与网格加密分别给出这两个误差的量级
\end{assumptionlist}

\subsection{假设代价的汇总}

上述 15 条假设可分为三类。第一类是\textbf{几何与守恒的简化}（A1$\sim$A6、A10），它们决定模型能描述哪些物理过程；其中 A1 与 A10 的代价可以用题目数据做量级估计（端面占比 $0.08$、体积收缩比例），A5 的影响则无法用本题数据量化。第二类是\textbf{闭合关系的形式}（A7$\sim$A9、A11），它们决定边界与物性的写法；其中 A7 是唯一``形式由题目给出、但物理基准未经证明''的一条，本文把它列为首要结构性风险。第三类是\textbf{数据与外推}（B1$\sim$B4），它们决定输入数据的可信区间；其中 B2 覆盖了 $93\%$ 的时长，是最大的不确定性来源。

需要特别说明的是，这三类误差\textbf{不能靠加密网格或收紧时间容差消除}：加密只能降低离散误差，而 $10^{-15}$ 量级的守恒残差也只证明``离散格式的收支是闭合的''，并不证明物理闭合正确。第~\ref{sec:validation} 节据此把验证结果分为``数值可信度''与``物理结构性局限''两部分分别陈述。
